\documentclass[]{article}

\usepackage[utf8]{inputenc}
\usepackage[english]{babel}
\usepackage{amsmath}

%For References and quotes
\usepackage[
backend=biber,
style=apa,
citestyle=apa,
sorting=nyt
]{biblatex}
\usepackage{csquotes}
\usepackage{hyphenat}

%For Links
\usepackage{hyperref}

%For Images
\usepackage{graphicx}
\graphicspath{ {./images/} }

%For Swot Analysis Matrixes
\usepackage{times}
\usepackage{tikz}
\usetikzlibrary{matrix}

%For TheoreticalFramework
\usepackage{varwidth}
\usepackage{longtable}
\newcommand{\textSR}{Server-side rendering}
\newcommand{\textSECO}{Server-side encoding}
\newcommand{\textSECR}{Server-side encrypting}
\newcommand{\textN}{Networking}
\newcommand{\textCDCO}{Client-side decoding}
\newcommand{\textCDCR}{Client-side decrypting}
\newcommand{\textCRD}{Client-side rendering / displaying}

\usepackage[acronym]{glossaries}

\makenoidxglossaries

\newacronym{gpu}{GPU}{Graphics Processing Unit}
\newacronym{cpu}{CPU}{Central Processing Unit}
\newacronym{ux}{UX}{User Experience}
\newacronym{vr}{VR}{Virtual Reality}
\newacronym{ms}{ms}{Milliseconds}
\newacronym{mtp}{MTP}{Motion\hyp{}to\hyp{}Photon}
\newacronym{hmd}{HMD}{Head Mounted Display}
\newacronym{fov}{FoV}{Field\hyp{}of\hyp{}View}
\newacronym{rgb}{RGB}{Red Green Blue}
\newacronym{qoe}{QoE}{Quality of Experience}
\newacronym{irh}{IRH}{Industrial Reality Hub}
\newacronym{aws}{AWS}{Amazon Web Services}
\newacronym{iam}{IAM}{Identification and Access Management}
\newacronym{tls}{TLS}{Transport Layer Security}
\newacronym{aes}{AES}{Advanced Encryption Standard}
\newacronym{fps}{FPS}{Frames per second}
\newacronym{vm}{VM}{Virtual Machine}
\addbibresource{references.bib}


\begin{document}

\begin{titlepage}
   \begin{center}
       \vspace*{1cm}

       \textbf{Final Implementation Plan}

       \vspace{0.5cm}
       Cloud VR. Secure, Fast and Distributed Virtual Reality Solutions.
       \vspace{1.5cm}

       \textbf{Leon Koster}

       \vfill
            
       \vspace{0.8cm}
     
      \includegraphics[width=0.5\textwidth]{university}
            
       Academie Creative Technology (ACT)\\
       Saxion University of Applied Sciences\\
       Netherlands\\
       May 2020
            
   \end{center}
\end{titlepage}

\tableofcontents
\newpage
\printnoidxglossary[type=\acronymtype]

\section{Introduction}
\label{sec:int}
Recent developments in the field of \acrfull{vr} offer all kinds of opportunities in the field of training
and entertainment. For training purposes, the audiovisual entry into a virtual world is where the biggest value is.
The capabilities of artificial environments allow users to manage scenarios and 
experiences that cannot be simulated in the real world. \acrshort{vr} also allows users to access the virtual training at any time
and less physical facilities are required for exercises. 
Example \acrshort{vr} experiences include training maintenance at high altitudes (such as windmills), working under heavy loads and weather conditions in construction (Strukton) or maintenance on naval ships (Thales). These companies (and more) form the \acrfull{irh}, which is one of the stakeholders of this project.The \acrshort{irh} is an industry consortium of 17 partners in AR/VR - The European digital innovation hub for industrial applied Augmented and Virtual Reality, stimulating cooperation and innovation between companies, government and knowledge institutes, resulting in world class business, knowledge and facilities. The hub is the AR/VR Fieldlab in the Dutch Smart Industry program and is recognized by the European commission as Digital Innovation Hub for industrial AR/VR \parencite{irh}.

\newpage
\section{Reason for Assignment}
\subsection{Graduation Assignment}
The assignment for this graduation is to analyse the current state of cloud \acrshort{vr} streaming. This includes existing commercial services, open source libraries/frameworks, relevant research papers and experimental technologies. The 4 research directions for this project are system architecture for a cloud \acrshort{vr} system, Latency, Multi-User Experiences and \acrshort{gpu} Scaling. This paper will focus on researching the system architecture and latency improvements. These findings will be presented in the form of a theoretical framework and a literature review.

Armed with the information from the initial part of the research, the student and the stakeholders will agree on a suitable research focus for the reminder of the projec, including a proof of concept focused on reducing latency.
\subsection{Client outline}
The stakeholders in this project are the student (Leon Koster), Saxion's XR Lab (Matthijs van Veen, Yiwei Jiang) and the \acrfull{irh}.

\newpage
\section{Objectives of the client}
Together with the companies from the Industrial Reality Hub mentioned in the Client outline, Saxion wants to investigate how virtual reality can be rendered in the cloud in a safe and efficient manner. This involves looking at state-of-the art technology in the field of virtual reality, cloud computing, rendering and machine learning for one complete Cloud\acrshort{vr} pipeline. 

In the bigger picture this project is a multi phase project, with this project being the initial phase. It aims to give an overview of the current situation regarding cloud \acrshort{vr} streaming and focus on system architecture and latency optimization. These informations will subsequently be utilized in the later stages of the over arching research project.

\section{Preliminary Problem Statement}
One of the essentials for a good VR experience is a powerful computer system to render semi-realistic
worlds. However, there are two problems here. First, this type of system is not available in every location.
Certainly if realistic images have to be rendered in the simulation, it requires specialized and expensive machines that are difficult to move. 

The second problem is that for rendering the VR training scenario, all kinds of data about the scenarios need to be available on the system. 
This can pose a problem when it concerns sensitive information, for example about all kinds of information
defense systems or business sensitive information.

Focusing on these problems will lay the foundation for future research, to make CloudVR streaming a mature technology.

The aim of this project is to investigate the feasability of a streaming based VR approach with current cutting edge technology.
Qualitative research methods will be used to gain in-depth insights about existing solutions and the current state of research into this topic.
The data will be contextualized via a literature review of recent research papers and an analysis of the capabilites of existing solutions when applied to the research problem.
\section{Theoretical Framework}

In order to thoroughly understand the aim and subject of the research, it is important to explore different existing solutions and literature. Therefore, the subjects that will be discussed in the  following theoretical framework are Cloud Streaming/Cloud Computing and Virtual Reality. Within this theoretical framework, definitions of the subjects will be given as well as current insights into these subjects. The topics reflect knowledge needed to understand the problem space. Together all of the topics make up the 360 scan.

\subsection{Cloud Streaming/Cloud Computing}

\subsubsection{Definition}
According to Armbrust et al. (2010) Cloud computing is defined as follows: 
"Cloud computing refers to both the applications delivered as services over the Internet and the hardware and systems software in the data centers that provide those services." \parencite[]{aviewoncc}
We can then further define Cloud streaming as the applications that are delivered over the internet as a service.

\subsubsection{Existing Solutions and Technology}
Several commercial gaming Cloud streaming services already exist, such as Google Stadia \parencite{stadia}, XBox XCloud \parencite{xcloud} and Nvidia GeForceNow \parencite{geforcenow}. There is also a variety of Infrastructure\hyp{}as\hyp{}a\hyp{}service (IaaS) platforms, such as Amazon's AWS \parencite{aws}, Microsoft's Azure \parencite{azure} and Google's Cloud Platform \parencite{gcp}. For more information about these applications and technologies, please refer to the \hyperref[sec:lit]{Literature Analysis}.

\subsubsection{System Architecture (for a cloud system)}


\subsubsection{Latency}

\subsection{Virtual Reality}

\subsubsection{Definition}
Early works, such as the 1996 article "Virtual Reality in Scientific Visualizations"  or the 2003 book "Virtual Reality Technology", define Virtual Reality as: "The use of computers and human\hyp{}computer interfaces to create the effect of a three\hyp{}dimensional world containing interactive objects with a strong sense of three\hyp{}dimensional presence."\parencite{vrsv} or "A Simulation in which computer graphics is used to create a realistic looking world. Moreover, the synthetic world is not static, but responds to the users input. This defines a key feature of Virtual Reality, \textit{real\hyp{}time interactivity}" \parencite{vrtech}.

\subsubsection{Physical Constraints}

As mentioned before, when developing a \acrshort{vr} application, there are a few physical constraints that developers need to be aware of. The most important threshold to know is the 20 \acrshort{ms} \acrfull{mtp} delay. Upon input from the \acrfull{hmd}, the developer has to display a new rendered image within an average of 20 \acrshort{ms} to avoid motion sickness for users. The more this threshold can be undercut, the better the chances to have an acceptable gameplay experience without motion sickness. Interaction input, such as the input from the controllers, can safely be processed at delays of $>$100 \acrshort{ms} without any negative repercussions in terms of \acrfull{ux}. For more information, please refer to the \hyperref[sec:lit]{Literature Analysis}.
\section{Final Problem Statement}
\section{Research Questions}

\subsection*{Main Question 1:} What is the current state of cloud \acrfull{vr} streaming?
\paragraph{Sub Question 1:} How (in)effective are existing solutions when applied to the cloud streaming \acrshort{vr} context?
\paragraph{Sub Question 2:} What research has been done on the shortcomings from question 1 ? Which area still requires the most research ?
\paragraph{Sub Question 3:} What are the most important considerations when designing an architecture for cloud \acrshort{vr} streaming?
\paragraph{Sub Question 4:} What methods are the most efficient way to reduce \acrfull{mtp} Latency to $\leq$20\acrshort{ms}?


\section{Approach}
To answer the research questions, I will create a prototype cloud \acrshort{vr} system. I aim to include as many pre-existing solutions as possible in order to accelerate the prototyping speed. Initially I want to design and implement a Minimal Viable Product with components I understand. If time allows I plan to create more prototypes with different technologies, in order to be able to compare the effectiveness of different solutions when applied to the research problem.

As I outlined earlier, the main focus for this report is on latency reduction and security. To compare solutions and understand their effectiveness I will measure \acrfull{mtp} latency as the primary benchmark. It is harder to measure security, but by understanding what the state-of-the-art security practices are and implementing them to the best of my ability, I aim to create a solution that keeps the clients data safe.

To obtain the needed knowledge to execute upon the presented approach, I plan to refer to Research Papers and Experts. There is a multitude of research papers for individual parts of the prototype (Networking, Latency Reduction, Encryption, Encoding, etc.) that contain concentrated knowledge. To assert which information are useful when applied to the research problem and to discover potential pitfalls, I plan on talking to experts that have real world expertise in developing similar products. I want to source experts from inside the \acrfull{irh} and outside of it (such as the ZeroLight developers \parencite{zerolight5g}) in order to maximise my chances of picking the right information/technology. 

\newpage
\section{Scope}
In this project/report I will outline the available technology that is useful for a cloud \acrshort{vr} system. I will furthermore work to create a Prototype solution that demonstrates the viability of cloud \acrshort{vr}, with a focus on latency reduction and security.

I will not create a production ready product and I will not be responsible for setting up a networking system.

\newpage
\section{Graduation Schedule}
I was asked to give an overview of the time line for this project. These times are from the graduation manual, translated to the time frame of my graduation project. 
\paragraph{Start of Project}
The project started on the 4th May 2020.
\paragraph{Preliminary Implementation Plan}
To be handed in at the end of week 2: 17th May 2020 23:59
\paragraph{Final Implementation Plan}
To be handed in at the end of week 5: 7th June 2020 23:59
\paragraph{Graduation Report Draft}
To be handed in at the end of week 10: $\sim$12th July 2020 23:59
\paragraph{Final Graduation Report}
To be handed in at the end of week 16: $\sim$23rd August 2020 23:59

\newpage
\printbibliography
\end{document}